\documentclass[11pt,a4paper]{article}
\usepackage[utf8]{inputenc}
\usepackage[T1]{fontenc}
\usepackage[spanish]{babel}
\usepackage{hyperref}
\usepackage{geometry}
\geometry{margin=2.5cm}
\usepackage{parskip}

\title{331 lo que firma lo que muta tiene que estar fuera}
\author{Alejandro Valenzuela Rivera\\\small Universidad Autónoma de Nuevo León}
\date{\today}
\begin{document}
\maketitle

\begin{quote}
\textbf{Sesgo declarado.} Este trabajo se sostiene desde el sesgo declarado de su autor, y no finge neutralidad. Lo que afirma se puede tumbar, y ese es el punto: el falsificador está escrito, no escondido.
\end{quote}

\section*{Resumen}
Si lo que muta se firma a sí mismo, no hay borde: hay un alma renombrada al revés.

\section{Por qué se sostiene}
Un borde que se firma desde adentro no es un borde: es una etiqueta.

\section{De dónde viene}
La pregunta de la autoridad: si el RUTA muta, \textquestiondown{}qué lo firma?

\section{El caso}
La Source Wall: el borde que solo sirve si lo firma algo que no es el que lo cruza.

\section{Método}
\textquestiondown{}Hay alguien fuera de lo que mutó que lo pueda firmar?

\section{Límites, y qué lo tumba}
Un caso donde el que muta pueda firmarse sin perder el borde.

\section{Lo que queda abierto}
\textquestiondown{}La firma de lo que mutó viene de afuera, o de lo mismo que mutó?

\section*{Cómo se leyó esta tesis}
Esto se escribió en el repositorio man-made-empathy, del que este candado es una pieza: candados/06 del caso/331 lo que firma lo que muta tiene que estar fuera.md. La máquina que lo compuso lo hizo desde este mismo archivo, y no agregó nada que no estuviera escrito. Si algo de esto es falso, se dice en el mismo sitio, y se corrige con fecha.

\end{document}
